%!TEX root = ../main.tex
\chapter*{Preface}\label{pref}
\addcontentsline{toc}{chapter}{Preface}

\lettrine[lines=2, lhang=0.33, loversize=0.1]{I}{ have}
always considered myself a problem solver. Whether I am writing code,
working through a chemical reaction, or trying to understand why
something does not work as expected, I am drawn to problems that require
patience, creativity, and a willingness to approach things differently.
This is ultimately what attracted me to both chemistry and computer
science. Chemistry has taught me to think beyond what I can immediately
observe and to consider the countless interactions that explain the
world around us. Computer science has taught me that failure is often
part of progress, and that solving a problem sometimes requires starting
over, questioning an assumption, or finding an entirely different
approach. Together, these disciplines have shaped how I think, how I
learn, and the kind of work I hope to pursue.

The inspiration for this Independent Study came from a conversation
with Dr.~Paul Edmiston, who introduced me to a challenge involving
molecular-networking software. His students were exploring how
molecular networking could help identify chemically modified substances
relevant to forensic toxicology and support investigations involving
drug-facilitated crimes. However, understanding the computational
processes behind the software and evaluating the reliability of its
results presented a challenge. Dr.~Edmiston needed someone willing to
look beneath the surface of the software and understand how it worked.
As a brother to seven sisters, the opportunity to contribute to research
with potential implications for personal safety and justice resonated
with me deeply. Beyond its scientific importance, I recognized a chance
to apply my knowledge of chemistry and computer science to a problem
that could make a meaningful difference in people's lives.

What made this opportunity especially exciting was that it brought
together the two fields I had spent my undergraduate career studying.
I have never viewed chemistry and computer science as competing
interests. Instead, I see them as disciplines that strengthen one
another. Chemistry presents complicated questions about the world,
while computer science provides tools for investigating those questions
at a scale and level of detail that would otherwise be difficult to
achieve. Although the research eventually developed into a project
centered on molecular networking and the identification of unknown
features in remote environmental samples, that initial conversation
remained central to my motivation: I wanted to understand how the
software worked and how it could be made more useful and reliable for
the scientists who depend on it.

My experiences outside the classroom have further shaped that goal.
During my internship as a software developer at the Federal Reserve
Bank of St.~Louis, I learned the importance of writing dependable
software, testing my work, and understanding the consequences of
decisions made within complex systems. Working with cloud infrastructure
and software architecture has shown me how careful design affects
reliability and long-term usefulness. As a coding instructor, I have
taught hundreds of students how to approach difficult programming
problems, often reminding them of the same lesson I continue to learn
myself: not understanding something immediately does not mean it cannot
be understood. My previous computational chemistry work, including
developing tools to explore multistep organic synthesis using large
reaction datasets, also helped me recognize how powerful software can
become when it is built with a clear understanding of the science it
supports.

After graduation, I hope to pursue a career developing software for
chemists and other scientific researchers. I want to build tools that
do more than produce results. I want researchers to understand those
results, explore the evidence behind them, and have confidence in the
computational processes they use. This Independent Study represents my
first substantial attempt to bring that ambition into practice. It has
also challenged me to recognize that being both a chemist and a computer
scientist does not mean having every answer. Instead, it means being
willing to learn from both disciplines, recognize the limitations of
each, and do the work necessary to connect them.

I would not have reached this point without the faculty members who
have challenged, supported, and believed in me throughout my time at
The College of Wooster.

Within the Computer Science Department, I want to thank Dr.~Heather
Guarnera for consistently helping me find ways to make what once seemed
impossible feel achievable. Your encouragement has pushed me to think
bigger about what I am capable of accomplishing. I thank Dr.~Drew
Guarnera, my first computer science professor, for helping establish the
foundation upon which so much of my work now rests. I am also grateful
to Dr.~Palmer for contributing to my growth as a computer scientist and
to Dr.~Musgrave, my Senior Independent Study advisor, for guiding me
through the computational challenges of this interdisciplinary project.

Within the Chemistry Department, I owe a special thank-you to Dr.~Faust,
my Senior Independent Study advisor and one of my greatest champions.
Your belief in my abilities, your encouragement, and your willingness
to help me grow have meant more than I can adequately express. I thank
Dr.~Martin, my academic advisor, for helping guide my academic journey
and supporting my interests across two demanding disciplines. Dr.~Snider,
thank you for the difficult but necessary conversation that helped me
recognize where I needed to grow. Sometimes the advice we need most is
not the advice we expect to hear. I thank Dr.~Bonvallet for being a
lecturer whose teaching has consistently made learning chemistry
engaging. I thank Dr.~Gallo for helping me fall in love with chemistry
and recognize just how much there is to discover within it. Dr.~Baker,
thank you for being one of my favorite professors and for the simple
kindness of always being someone I look forward to greeting. Finally,
Dr.~Feierabend, thank you for continuing to challenge me to pursue
excellence, particularly during moments when I have struggled to
recognize my own potential.

I also want to thank Dr.~Edmiston for presenting the problem that
ultimately helped me discover the direction I want to take my career.
Sometimes the most important opportunities begin with a conversation
and a question that nobody has fully answered.

This Independent Study is more than a requirement for graduation. It
represents the intersection of my interests, my experiences, and the
people who have helped shape my education. More importantly, it is a
beginning. I entered college wanting to understand how things work.
I hope to leave having taken my first meaningful steps toward building
things that help others understand the world.